% from ega4-17-fr.tex
\oldpage[IV-4]{108}

\begin{corollary}[17.16.5]
\label{IV.17.16.5}
Soit $f:X\to S$ un morphisme surjectif localement de présentation finie. Alors il existe un morphisme $g:S'\to S$ surjectif, localement de présentation finie et localement quasi-fini, tel qu'il existe un $S$-morphisme $S'\to X$ (autrement dit tel qu'il existe une $S'$-section de $X'=X\times_S S'$). Si $S$ est quasi-compact (resp. quasi-compact et quasi-séparé), on peut supposer $S'$ affine (resp. $S'$ affine et le morphisme $g$ de présentation finie et quasi-fini).
\end{corollary}

Il suffit de prouver le corollaire en supposant $S$ affine: $S$ est en effet réunion d'une famille $(S_\alpha)$ d'ouverts affines, et si pour chaque $\alpha$, $S'_\alpha$ est affine et le morphisme $g_\alpha:S'_\alpha\to S_\alpha$ répond à la question et est de présentation finie et quasi-fini, alors en prenant pour $S'$ le préschéma somme des $S'_\alpha$, $g:S'\to S$ coïncidant avec $g_\alpha$ dans chacun des $S'_\alpha$, ce morphisme répond à la question; en outre, si $S$ est quasi-compact, on peut supposer la famille $(S_\alpha)$ finie, donc $S'$ affine; si de plus $S$ est quasi-séparé, les immersions $S_\alpha\to S$ sont de présentation finie (I, 6.2), donc il en est de même de $g$.

Considérons donc le cas où $S$ est affine: on forme alors les morphismes finis et de présentation finie $S_i''\to S_i$ de (17.16.4); comme les immersions $S_i\to S$ sont de présentation finie (les $S_i$ étant affines), les morphismes $g_i'':S_i''\to S_i\to S$ sont de présentation finie et quasi-finis, et on répond à la question en prenant $S'$ égal à la somme des $S_i''$ et $g$ coïncidant avec $g_i''$ dans chacun des $S_i''$.

\begin{proposition}[17.16.6]
\label{IV.17.16.6}
Soient $S$ un préschéma quasi-compact et quasi-séparé, $f:X\to S$ un morphisme de présentation finie tel que, pour tout $s\in S$, $X_s=f^{-1}(s)$ soit un préschéma propre sur $k(s)$. Alors il existe une famille finie $(S_i)_{i\in I}$ de sous-préschémas affines de $S$, de présentation finie sur $S$, deux à deux disjoints, de réunion $S$ et tels que, pour tout $i$, le morphisme $X_i=X\times_S S_i\to S_i$ soit propre et plat. Si de plus, pour tout $s\in S$, $X_s$ est fini sur $k(s)$, on peut prendre les $S_i$ tels que chacun des morphismes $X_i\to S_i$ se factorise en
\[
  X_i\xrightarrow{u_i}X_i'\xrightarrow{h_i}S_i
\]
où $u_i$ et $h_i$ sont finis et localement libres (18.2.7), $h_i$ est étale, $u_i$ est radiciel et surjectif.
\end{proposition}

On suit une marche analogue à celle de (17.16.4). On se ramène d'abord au cas où $S$ est affine et où $f$ est plat, en utilisant le théorème de platitude générique (8.9.5) (on observera que les sous-préschémas $S_i'$ définis dans la démonstration de (8.9.5) sont affines). Puisque $f$ est de présentation finie, on peut ensuite écrire $X=X_0\times_{S_0}S$, où $S_0$ est noethérien, $X_0\to S_0$ un morphisme de type fini et plat (11.2.7); en outre chaque fibre $(X_0)_{s_0}$ est encore propre sur $k(s_0)$, comme il résulte de (2.7.1, (vii)) et on est donc ramené au cas où $S$ est noethérien. Par récurrence noethérienne et application du procédé de (8.1.2, $a$)) (en utilisant aussi cette fois (8.10.5, (xii))), on est finalement ramené, pour prouver la première assertion, au cas où $S$ est le spectre d'un corps, qui résulte trivialement de l'hypothèse (avec $S_i=S$). On traite de même le cas où $X_s$ est supposé fini sur $k(s)$ pour tout $s\in S$ (il faut cette fois utiliser (2.7.1, (xv)), (8.10.5, (x), (iv), (vi) et (vii)), (17.7.8) et (2.1.12)), et on est ramené à prouver la dernière assertion de l'énoncé lorsque $S=\Spec(k)$ est le spectre d'un corps. Mais alors $X=\Spec(A)$, où $A$ est une $k$-algèbre finie (I, 6.4.4); donc $A$ est composée directe de $k$-algèbres finies $A_j$ qui sont des anneaux locaux $(1\leq j\leq m)$; puisque $A_j$ est un anneau artinien et une

